\chapter{Hiện thực - User Portal}

\section{Tổng quan User Portal}

User Portal là ứng dụng di động chính dành cho người dùng cuối (bệnh nhân), cung cấp các tính năng:
\begin{itemize}
    \item Tìm kiếm và đặt lịch hẹn dịch vụ y tế/làm đẹp
    \item Tư vấn trực tiếp với trợ lý AI (Dr. AI Assistant)
    \item Quản lý lịch hẹn và lịch sử giao dịch
    \item Thanh toán qua nhiều phương thức (Stripe, Momo, e-wallet)
    \item Giao tiếp với nhà cung cấp dịch vụ
\end{itemize}

\section{Kiến trúc giao diện User Portal}

\subsection{Màn hình chính (Home)}

Màn hình chính là điểm khởi đầu của người dùng, cung cấp truy cập nhanh đến các tính năng chính:

\begin{figure}[H]
\centering
\begin{subfigure}[b]{0.48\linewidth}
    \centering
    \includegraphics[width=\linewidth]{Images/Hiện\ thực/user/home_light.png}
    \caption{Giao diện chế độ sáng với Quick Actions và Recommend For You}
\end{subfigure}
\hfill
\begin{subfigure}[b]{0.48\linewidth}
    \centering
    \includegraphics[width=\linewidth]{Images/Hiện\ thực/user/home_service_detail.png}
    \caption{Chi tiết dịch vụ được đề xuất}
\end{subfigure}
\caption{Màn hình Home - Trung tâm của User Portal}
\label{fig:user_home}
\end{figure}

\textbf{Thành phần chính:}

\begin{itemize}
    \item \textbf{User Profile Header:} Hiển thị tên người dùng, icon settings, và giỏ hàng. Người dùng có thể nhấp vào header để xem thông tin tài khoản.
    
    \item \textbf{Quick Actions:} Hai nút bấm chính với icon hấp dẫn:
    \begin{itemize}
        \item \textit{Book Appointment}: Dẫn người dùng tới màn hình chọn chuyên gia và lập lịch hẹn
        \item \textit{AI Health Assistant}: Mở chat với Dr. AI để nhận tư vấn sức khỏe
    \end{itemize}
    
    \item \textbf{Recommend For You:} Danh sách dạo động (carousel) các dịch vụ được khuyến nghị dựa trên:
    \begin{itemize}
        \item Lịch sử đặt lịch trước đó của người dùng
        \item Đánh giá dịch vụ cao
        \item Dịch vụ mới nhất từ các nhà cung cấp
    \end{itemize}
    
    \item \textbf{Tỷ lệ đánh giá:} Mỗi dịch vụ hiển thị sao (★) và số lượt đánh giá để người dùng có cơ sở quyết định.
    
    \item \textbf{Bottom Navigation:} Thanh điều hướng 5 tab: Home, Orders, Chat, Notifications, Profile
\end{itemize}

\textbf{Tính năng kỹ thuật:}

\begin{itemize}
    \item \textbf{State Management:} Sử dụng BLoC pattern để quản lý trạng thái danh sách dịch vụ recommended. Khi người dùng mở lại Home, dữ liệu được cache để tránh load lại không cần thiết.
    
    \item \textbf{Image Loading:} Hình ảnh dịch vụ được tải với fade-in animation và placeholder, giảm thiểu cảm giác ``trắng'' khi mạng chậm.
    
    \item \textbf{Pull-to-Refresh:} Người dùng có thể kéo xuống để làm mới danh sách dịch vụ mà không cần nhấn vào ô tìm kiếm.
\end{itemize}

\subsection{Quản lý lịch hẹn (Appointments)}

Màn hình Appointments cho phép người dùng xem các lịch hẹn được phân loại thành bốn tab:

\begin{figure}[H]
\centering
\begin{subfigure}[b]{0.48\linewidth}
    \centering
    \includegraphics[width=\linewidth]{Images/Hiện\ thực/user/appointments_upcoming.png}
    \caption{Tab Upcoming - Lịch hẹn sắp tới}
\end{subfigure}
\hfill
\begin{subfigure}[b]{0.48\linewidth}
    \centering
    \includegraphics[width=\linewidth]{Images/Hiện\ thực/user/appointments_dental.png}
    \caption{Danh sách dịch vụ từ Healytics Dental}
\end{subfigure}
\caption{Màn hình Appointments - Quản lý lịch hẹn}
\label{fig:user_appointments}
\end{figure}

\textbf{Cấu trúc tab:}

\begin{itemize}
    \item \textbf{Pending:} Lịch hẹn đang chờ xác nhận từ nhà cung cấp. Người dùng có thể hủy hoặc chỉnh sửa chi tiết.
    
    \item \textbf{Upcoming:} Lịch hẹn đã được xác nhận và sắp diễn ra trong tương lai. Mỗi lịch hẹn hiển thị:
    \begin{itemize}
        \item Tên dịch vụ (ví dụ: ``Full Body Massage 60 Min'')
        \item Chuyên gia phụ trách (ví dụ: ``Sarah Mitchell'')
        \item Thời gian check-in (ví dụ: ``Upcoming'', ``Check in at 08:00 AM'')
        \item Địa chỉ và vị trí trên bản đồ
        \item Nút ``Details'' để xem chi tiết
    \end{itemize}
    
    \item \textbf{Completed:} Lịch hẹn đã hoàn thành. Người dùng có thể xem lại thông tin và để lại đánh giá.
    
    \item \textbf{Canceled:} Lịch hẹn đã bị hủy (bởi người dùng hoặc nhà cung cấp). Hiển thị lý do hủy nếu có.
\end{itemize}

\textbf{Tính năng kỹ thuật:}

\begin{itemize}
    \item \textbf{Real-time Status Sync:} Trạng thái của mỗi lịch hẹn được cập nhật real-time thông qua WebSocket khi nhà cung cấp xác nhận hoặc đánh dấu hoàn thành.
    
    \item \textbf{Push Notification:} Khi lịch hẹn gần tới (ví dụ 1 giờ trước), hệ thống gửi reminder notification.
    
    \item \textbf{Lazy Loading:} Danh sách lịch hẹn được tải từng batch (10 item) khi người dùng scroll xuống, tránh tải toàn bộ.
\end{itemize}

\subsection{Chọn chuyên gia và lịch (Select Specialist)}

Quy trình đặt lịch gồm ba bước, bắt đầu bằng việc chọn chuyên gia phù hợp và slot thời gian:

\begin{figure}[H]
\centering
\begin{subfigure}[b]{0.48\linewidth}
    \centering
    \includegraphics[width=\linewidth]{Images/Hiện\ thực/user/select_specialist.png}
    \caption{Chọn chuyên gia, ngày và giờ}
\end{subfigure}
\hfill
\begin{subfigure}[b]{0.48\linewidth}
    \centering
    \includegraphics[width=\linewidth]{Images/Hiện\ thực/user/available_time_slots.png}
    \caption{Hiển thị các slot thời gian khả dụng}
\end{subfigure}
\caption{Màn hình Select Specialist - Chọn chuyên gia và thời gian}
\label{fig:user_select_specialist}
\end{figure}

\textbf{Thành phần chính:}

\begin{itemize}
    \item \textbf{Select Specialist:} Danh sách các chuyên gia theo dạng card có hình ảnh avatar, tên, chuyên ngành. Card được chọn sẽ có viền xanh và checkmark.
    \begin{itemize}
        \item Mỗi chuyên gia hiển thị: Avatar (custom emoji/illustration), Họ tên, Chuyên ngành (ví dụ: ``Physical Therapy Technician'')
        \item Người dùng có thể scroll ngang hoặc nhấp ``See all'' để xem thêm
    \end{itemize}
    
    \item \textbf{Select Date:} Hiển thị 5 ngày liên tiếp (Hôm nay, Ngày mai, ...). Ngày được chọn hiển thị màu xanh.
    \begin{itemize}
        \item Format: Ngày (hôm nay, ngày tuần) và Số (4, 5, 6, 7, 8)
        \item Người dùng có thể scroll sang để xem thêm ngày
    \end{itemize}
    
    \item \textbf{Available Time Slots:} Phân chia thành ``Morning'' (biểu tượng mặt trời) và ``Afternoon'' (biểu tượng mặt trời). Mỗi slot hiển thị giờ (ví dụ: 09:00 AM, 09:30 AM, ...).
    \begin{itemize}
        \item Slot khả dụng: Viền xanh nhạt, text xanh
        \item Slot không khả dụng: Ẩn hoặc hiển thị với màu xám (disabled)
    \end{itemize}
    
    \item \textbf{Nút Continue:} Chuyển sang bước tiếp theo (Confirm Booking).
\end{itemize}

\textbf{Tính năng kỹ thuật:}

\begin{itemize}
    \item \textbf{Real-time Availability:} Danh sách slot khả dụng được lấy từ API backend, xét các yếu tố:
    \begin{itemize}
        \item Lịch làm việc của chuyên gia
        \item Các lịch hẹn đã được đặt
        \item Thời gian buffer giữa các appointment (ví dụ: 15 phút setup)
    \end{itemize}
    
    \item \textbf{Distributed Locking:} Khi người dùng chọn một slot, hệ thống sử dụng Redis Redlock để ngăn chặn race condition (nhiều user cùng chọn một slot).
    
    \item \textbf{Form State Persistence:} Nếu người dùng quay lại màn hình này, các lựa chọn trước đó (chuyên gia, ngày, giờ) được lưu lại trong state, tránh mất dữ liệu.
\end{itemize}

\subsection{Xác nhận đặt lịch (Confirm Booking)}

Màn hình Confirm Booking (Step 3 of 3) hiển thị toàn bộ chi tiết đơn hẹn và cho phép thanh toán:

\begin{figure}[H]
\centering
\begin{subfigure}[b]{0.48\linewidth}
    \centering
    \includegraphics[width=\linewidth]{Images/Hiện\ thực/user/confirm_booking_1.png}
    \caption{Chi tiết dịch vụ và chuyên gia}
\end{subfigure}
\hfill
\begin{subfigure}[b]{0.48\linewidth}
    \centering
    \includegraphics[width=\linewidth]{Images/Hiện\ thực/user/confirm_booking_2.png}
    \caption{Chi phí, địa chỉ và thanh toán}
\end{subfigure}
\caption{Màn hình Confirm Booking - Xác nhận chi tiết}
\label{fig:user_confirm_booking}
\end{figure}

\textbf{Thành phần chính:}

\begin{itemize}
    \item \textbf{Step Indicator:} Hiển thị ``Step 3 of 3'' ở phía trên cùng, giúp người dùng biết đã bao xa trong quy trình.
    
    \item \textbf{Specialist Info:} Avatar, tên chuyên gia (ví dụ: ``Sarah Mitchell''), chuyên ngành (``Physical Therapy Technician''), với badge xác thực (checkmark).
    
    \item \textbf{Appointment Details:}
    \begin{itemize}
        \item Ngày và giờ: ``May 5, 2026'' ``11:00 AM''
        \item Danh mục và dịch vụ: ``CATEGORY \& SERVICE'' với ảnh thumbnail, tên dịch vụ, mô tả ngắn
    \end{itemize}
    
    \item \textbf{Location Details:} Tên cơ sở (``MindSkin Clinic''), địa chỉ đầy đủ, nút bản đồ để xem vị trí.
    
    \item \textbf{Price Breakdown:}
    \begin{itemize}
        \item Subtotal: Giá gốc (ví dụ: 350,000 VND)
        \item Service Fee: Phí dịch vụ (ví dụ: -51,000 VND, nếu có khuyến mãi)
        \item Total Amount: Tổng cuối cùng (ví dụ: 299,000 VND)
    \end{itemize}
    
    \item \textbf{Nút Confirm \& Pay:} Nút chính dẫn tới màn hình thanh toán.
\end{itemize}

\textbf{Tính năng kỹ thuật:}

\begin{itemize}
    \item \textbf{Price Calculation Engine:} Backend tính toán giá cuối cùng xét:
    \begin{itemize}
        \item Giá cơ bản của dịch vụ
        \item Khuyến mãi hiện tại (voucher, sale)
        \item Coin/point của người dùng (nếu dùng)
        \item Thuế (VAT) nếu có
    \end{itemize}
    
    \item \textbf{Order Draft Creation:} Hệ thống tạo một ``draft order'' tạm thời với TTL (Time To Live) 10 phút. Nếu quá thời gian mà chưa thanh toán, draft này sẽ hết hạn và slot thời gian sẽ được giải phóng.
\end{itemize}

\subsection{Trò chuyện với trợ lý AI (Chat with Dr. AI)}

Tính năng AI Assistant cung cấp tư vấn sức khỏe 24/7 thông qua chat:

\begin{figure}[H]
\centering
\begin{subfigure}[b]{0.48\linewidth}
    \centering
    \includegraphics[width=\linewidth]{Images/Hiện\ thực/user/chat_history_ai.png}
    \caption{Lịch sử chat AI Session}
\end{subfigure}
\hfill
\begin{subfigure}[b]{0.48\linewidth}
    \centering
    \includegraphics[width=\linewidth]{Images/Hiện\ thực/user/ai_chat_conversation.png}
    \caption{Cuộc trò chuyện chi tiết}
\end{subfigure}
\caption{Màn hình Chat History - Tư vấn AI}
\label{fig:user_chat_history}
\end{figure}

\begin{figure}[H]
\centering
\begin{subfigure}[b]{0.48\linewidth}
    \centering
    \includegraphics[width=\linewidth]{Images/Hiện\ thực/user/ai_initial.png}
    \caption{Trợ lý AI ban đầu (greeting)}
\end{subfigure}
\hfill
\begin{subfigure}[b]{0.48\linewidth}
    \centering
    \includegraphics[width=\linewidth]{Images/Hiện\ thực/user/ai_conversation_detail.png}
    \caption{Chi tiết cuộc hội thoại với lời khuyên chuyên gia}
\end{subfigure}
\caption{Màn hình Dr. AI Assistant - Nội dung tư vấn}
\label{fig:user_ai_conversation}
\end{figure}

\textbf{Thành phần chính:}

\begin{itemize}
    \item \textbf{Chat History Tab:} Giao diện chuyên mục ``Chat History'' với hai tab chính:
    \begin{itemize}
        \item \textit{AI Session}: Lịch sử trò chuyện với Dr. AI Assistant
        \item \textit{Partner Chat}: Trò chuyện với các nhà cung cấp dịch vụ
    \end{itemize}
    
    \item \textbf{Conversation Listings:} Mỗi cuộc hội thoại hiển thị:
    \begin{itemize}
        \item Icon (suitcase-like logo)
        \item Tên cuộc hội thoại (ví dụ: ``SEED\_AI\_CONV\_DENTAL'', ``SEED\_AI\_CONV\_WELLNESS'')
        \item Thời gian tạo (ví dụ: ``8:13 AM'')
        \item Số hoặc trạng thái (ví dụ: mũi tên ``>'' để mở)
    \end{itemize}
    
    \item \textbf{Greeting Message:} Khi mở chat với AI lần đầu, trợ lý hiển thị:
    \begin{itemize}
        \item Icon Dr. AI đặc trưng (tròn với mặt nụ cười)
        \item Lời chào: ``Hi! I'm Dr. AI 👋''
        \item Hướng dẫn: ``Ask me about symptoms, appointments, or health questions.''
        \item Quick action buttons: ``Symptom Check'', ``Book Appointment'', và thêm
    \end{itemize}
    
    \item \textbf{Message Bubble:}
    \begin{itemize}
        \item AI Message (từ Dr. AI): Bubble màu xanh, đậm, có icon suitcase
        \item User Message (từ người dùng): Bubble xám nhạt, hiển thị bên phải
        \item Timestamp: Hiển thị giờ dưới mỗi tin nhắn
    \end{itemize}
    
    \item \textbf{Message Input:} Thanh input ở dưới cùng với:
    \begin{itemize}
        \item TextField để gõ tin nhắn
        \item Nút gửi (mũi tên trong vòng tròn xanh)
    \end{itemize}
    
    \item \textbf{Recommended Services:} Khi AI phát hiện vấn đề sức khỏe, nó gợi ý các dịch vụ phù hợp với giá và ``View →'' button để xem chi tiết.
\end{itemize}

\textbf{Tính năng kỹ thuật:}

\begin{itemize}
    \item \textbf{LLM Integration:} Chat AI được cung cấp bởi một Large Language Model (ví dụ: GPT-4, Llama). Các prompt được thiết kế để AI hoạt động như một y sĩ thực tế và gợi ý dịch vụ từ hệ thống.
    
    \item \textbf{Streaming Response:} Khi AI phản hồi, thay vì chờ toàn bộ tin nhắn hoàn thành, hệ thống stream từng token (từ) để cho cảm giác ``real-time'' và giảm thời gian chờ cảm nhận.
    
    \item \textbf{Context Preservation:} Hệ thống lưu trữ context của cuộc trò chuyện trên server. Nếu người dùng tắt ứng dụng và mở lại, toàn bộ lịch sử được khôi phục.
    
    \item \textbf{Service Recommendation Engine:} Khi người dùng mô tả triệu chứng, backend gọi AI service để phân tích và trả về danh sách dịch vụ phù hợp từ database, cùng với giải thích.
\end{itemize}

\subsection{Thanh toán (Checkout \& Payment)}

Quy trình thanh toán cho phép người dùng chọn phương thức thanh toán và xác nhận:

\begin{figure}[H]
\centering
\begin{subfigure}[b]{0.48\linewidth}
    \centering
    \includegraphics[width=\linewidth]{Images/Hiện\ thực/user/checkout_payment.png}
    \caption{Màn hình Checkout với các phương thức thanh toán}
\end{subfigure}
\hfill
\begin{subfigure}[b]{0.48\linewidth}
    \centering
    \includegraphics[width=\linewidth]{Images/Hiện\ thực/user/payment_confirmation.png}
    \caption{Xác nhận thanh toán thành công}
\end{subfigure}
\caption{Màn hình Checkout \& Payment Confirmation}
\label{fig:user_checkout}
\end{figure}

\textbf{Thành phần chính:}

\begin{itemize}
    \item \textbf{Order Items:} Hiển thị đơn hẹn được đặt:
    \begin{itemize}
        \item Tên cơ sở (``MindSkin Clinic'') với địa chỉ đầy đủ
        \item Hình ảnh thumbnail dịch vụ
        \item Tên dịch vụ, giờ đặt, chuyên gia
        \item Giá (với số lượng × nếu có)
    \end{itemize}
    
    \item \textbf{Redeem Coins:} Toggle cho phép người dùng dùng Coin/point hiện có để giảm giá. Hiển thị số coin hiện tại.
    
    \item \textbf{Payment Method:}
    \begin{itemize}
        \item ``Credit/Debit Card'' (VISA): Mặc định được chọn
        \item ``E-Wallet (Momo/ZaloPay)''
        \item ``Pay Later''
        \item Mỗi option hiển thị icon và radio button
    \end{itemize}
    
    \item \textbf{Total Amount:} Hiển thị tổng tiền cuối cùng (ví dụ: ``299,000đ'')
    
    \item \textbf{Nút Confirm Payment:} Xác nhận và tiến hành giao dịch với gateway thanh toán tương ứng.
    
    \item \textbf{Confirmation Modal:} Sau khi thanh toán thành công, hiển thị modal với:
    \begin{itemize}
        \item Icon checkmark xanh (thành công)
        \item Tiêu đề: ``Booking Confirmed!''
        \item Thông báo: ``Your booking has been created successfully.''
        \item Nút ``Done''
    \end{itemize}
\end{itemize}

\textbf{Tính năng kỹ thuật:}

\begin{itemize}
    \item \textbf{Payment Gateway Integration:} Hệ thống tích hợp với:
    \begin{itemize}
        \item \textbf{Stripe}: Cho thanh toán thẻ quốc tế
        \item \textbf{Momo}: Cho ví điện tử Việt Nam
        \item \textbf{PayLater}: Dịch vụ trả sau (nếu có partner)
    \end{itemize}
    
    \item \textbf{Webhook Handling:} Khi người dùng hoàn tất thanh toán ở gateway (ví dụ: Stripe redirect), server nhận webhook callback từ gateway để:
    \begin{itemize}
        \item Cập nhật trạng thái order thành ``paid''
        \item Gửi confirmation email
        \item Gửi push notification cho nhà cung cấp
    \end{itemize}
    
    \item \textbf{PCI Compliance:} Để tuân thủ PCI DSS, số thẻ không bao giờ được lưu trên server. Thay vào đó, frontend gửi token từ Stripe Tokenize API, backend chỉ nhận token này.
    
    \item \textbf{Idempotency:} Nếu người dùng nhấn Confirm Payment hai lần (do lag mạng), hệ thống sử dụng Idempotency Key để đảm bảo chỉ xử lý một lần, tránh double charge.
\end{itemize}

\subsection{Chi tiết dịch vụ (Service Detail)}

Khi người dùng chọn một dịch vụ, màn hình hiển thị toàn bộ thông tin:

\begin{figure}[H]
\centering
\begin{subfigure}[b]{0.48\linewidth}
    \centering
    \includegraphics[width=\linewidth]{Images/Hiện\ thực/user/service_detail_1.png}
    \caption{Thông tin cơ bản và đánh giá}
\end{subfigure}
\hfill
\begin{subfigure}[b]{0.48\linewidth}
    \centering
    \includegraphics[width=\linewidth]{Images/Hiện\ thực/user/service_detail_2.png}
    \caption{Mô tả dịch vụ, Gallery và Booking}
\end{subfigure}
\caption{Màn hình Service Detail - Chi tiết dịch vụ}
\label{fig:user_service_detail}
\end{figure}

\textbf{Thành phần chính:}

\begin{itemize}
    \item \textbf{Image Banner:} Ảnh lớn của dịch vụ (ví dụ: hình massage), có overlay gradient từ trên xuống, nút yêu thích (heart) ở trên phải, nút chia sẻ (share).
    
    \item \textbf{Service Category Tag:} Badge nhỏ với nền xanh nhạt, text xanh (ví dụ: ``RELAXATION MASSAGE'').
    
    \item \textbf{Service Title:} Tên dịch vụ (ví dụ: ``Full Body Massage 60 Min'')
    
    \item \textbf{Rating:} Sao (★) và số lượt review (ví dụ: ``4.5 (2 Reviews)''), kèm badge ``Verified Treatment''
    
    \item \textbf{Price:} Hiển thị giá per session (ví dụ: ``299.000đ / per session'')
    
    \item \textbf{About Treatment:} Phần mô tả chi tiết:
    \begin{itemize}
        \item Đoạn mô tả dài (collapsed mặc định, có nút ``Show more'')
        \item Mô tả chi tiết lợi ích, quy trình, v.v.
    \end{itemize}
    
    \item \textbf{Facility Tour:} Gallery hình ảnh cơ sở, có thể swipe ngang. Mỗi hình là ảnh quán café/phòng chờ/phòng massage.
    
    \item \textbf{Bottom Action:}
    \begin{itemize}
        \item Hiển thị giá cuối cùng ``Total Price: 299.000đ''
        \item Nút giỏ hàng để thêm vào giỏ
        \item Nút ``Select Specialist →'' để bắt đầu quy trình đặt lịch
    \end{itemize}
\end{itemize}

\textbf{Tính năng kỹ thuật:}

\begin{itemize}
    \item \textbf{Image Optimization:} Hình ảnh được optimize với nhiều kích thước:
    \begin{itemize}
        \item Thumbnail (100x100px) cho danh sách
        \item Medium (400x300px) cho màn hình chi tiết
        \item Full resolution cho gallery viewer
    \end{itemize}
    
    \item \textbf{Caching:} Thông tin dịch vụ được cache cục bộ (local storage) trong 24 giờ, giảm API request.
    
    \item \textbf{Related Services:} Ở phía dưới, hệ thống đề xuất các dịch vụ liên quan dựa trên cùng category hoặc từ cùng cơ sở.
\end{itemize}

\section{Kiến trúc dữ liệu và API}

\subsection{DTO (Data Transfer Object) chính}

\begin{verbatim}
// User API
UserDTO {
  id: UUID
  email: string
  firstName: string
  lastName: string
  phone: string
  avatarUrl?: string
  createdAt: timestamp
}

// Service DTO
ServiceDTO {
  id: UUID
  name: string
  description: string
  category: ServiceCategoryEnum
  price: decimal
  duration: integer (minutes)
  images: string[] (URLs)
  rating: decimal (0-5)
  reviewCount: integer
  provider: ProviderDTO
}

// Appointment DTO
AppointmentDTO {
  id: UUID
  userId: UUID
  serviceId: UUID
  specialistId: UUID
  appointmentAt: timestamp
  status: AppointmentStatusEnum (PENDING, CONFIRMED, COMPLETED, CANCELLED)
  totalPrice: decimal
  paymentStatus: PaymentStatusEnum
}

// Payment DTO
PaymentDTO {
  id: UUID
  appointmentId: UUID
  amount: decimal
  method: PaymentMethodEnum (CARD, MOMO, E_WALLET, PAY_LATER)
  gatewayReference: string
  status: PaymentStatusEnum (PENDING, SUCCESS, FAILED)
}

// Chat Message DTO
ChatMessageDTO {
  id: UUID
  conversationId: UUID
  senderType: SenderTypeEnum (USER, AI, PROVIDER)
  senderId: UUID
  content: string
  attachments?: Attachment[]
  createdAt: timestamp
}
\end{verbatim}

\subsection{API Endpoints chính}

\begin{table}[H]
\centering
\begin{tabularx}{\linewidth}{|l|X|X|}
\hline
\textbf{Endpoint} & \textbf{Phương thức} & \textbf{Mô tả} \\
\hline
\texttt{/v1/services} & GET & Lấy danh sách dịch vụ (có phân trang, lọc category) \\
\hline
\texttt{/v1/services/\{id\}} & GET & Chi tiết dịch vụ \\
\hline
\texttt{/v1/services/\{id\}/recommended} & GET & Dịch vụ được đề xuất \\
\hline
\texttt{/v1/appointments} & GET & Lấy danh sách lịch hẹn của user \\
\hline
\texttt{/v1/appointments} & POST & Tạo lịch hẹn mới (draft) \\
\hline
\texttt{/v1/appointments/\{id\}} & GET & Chi tiết lịch hẹn \\
\hline
\texttt{/v1/appointments/\{id\}/confirm} & POST & Xác nhận lịch hẹn sau thanh toán \\
\hline
\texttt{/v1/appointments/\{id\}/cancel} & POST & Hủy lịch hẹn \\
\hline
\texttt{/v1/services/\{id\}/availability} & GET & Lấy slot thời gian khả dụng \\
\hline
\texttt{/v1/payments} & POST & Khởi tạo thanh toán \\
\hline
\texttt{/v1/payments/\{id\}/confirm} & POST & Xác nhận thanh toán (từ webhook gateway) \\
\hline
\texttt{/v1/chat/conversations} & GET & Lịch sử cuộc trò chuyện \\
\hline
\texttt{/v1/chat/conversations} & POST & Tạo cuộc trò chuyện mới \\
\hline
\texttt{/v1/chat/messages} & POST & Gửi tin nhắn \\
\hline
\texttt{/v1/ai/chat} & POST & Gửi tin nhắn tới AI Assistant \\
\hline
\end{tabularx}
\caption{Danh sách API endpoints chính cho User Portal}
\label{tab:user_api_endpoints}
\end{table}

\section{Quy trình xử lý nghiệp vụ chính}

\subsection{Quy trình đặt lịch hẹn}

Toàn bộ quy trình đặt lịch từ lựa chọn dịch vụ đến xác nhận:

\begin{enumerate}
    \item \textbf{Người dùng chọn dịch vụ:} Trên màn hình Home hoặc danh sách dịch vụ, người dùng tap vào một dịch vụ.
    
    \item \textbf{Chuyển đến Select Specialist:} UI chuyển sang màn hình chọn chuyên gia, ngày, giờ.
    \begin{itemize}
        \item Backend gọi endpoint \texttt{/v1/services/\{id\}/availability} để lấy danh sách slot khả dụng.
        \item Danh sách này được lọc theo:
        \begin{itemize}
            \item Lịch làm việc của từng chuyên gia (kích hoạt từ Provider App)
            \item Các appointment đã được confirmed (trạng thái != CANCELLED)
            \item Buffer time (15 phút giữa các appointment)
        \end{itemize}
    \end{itemize}
    
    \item \textbf{Người dùng chọn chuyên gia, ngày, giờ:} Khi người dùng chọn xong, thông tin này được lưu trong state.
    \begin{itemize}
        \item Hệ thống gọi \texttt{POST /v1/appointments} để tạo draft order với TTL 10 phút.
        \item Draft order này không giữ lock trên database, chỉ là một bản ghi tạm thời.
    \end{itemize}
    
    \item \textbf{Chuyển đến Confirm Booking:} Màn hình hiển thị toàn bộ chi tiết:
    \begin{itemize}
        \item Backend tính toán giá cuối cùng (xét voucher, coin, VAT).
        \item Hiển thị tất cả thông tin chuẩn bị thanh toán.
    \end{itemize}
    
    \item \textbf{Người dùng chọn phương thức thanh toán:} Chuyển đến Checkout screen.
    
    \item \textbf{Thanh toán:}
    \begin{itemize}
        \item Nếu dùng **Credit Card**: Frontend gọi Stripe Tokenize API để nhận token, sau đó gửi token này tới backend.
        \item Nếu dùng **Momo**: Backend tạo QR code, người dùng scan và thanh toán ở Momo app.
        \item Backend gọi payment gateway \texttt{POST /v1/payments} để khởi tạo giao dịch.
    \end{itemize}
    
    \item \textbf{Webhook Callback:}
    \begin{itemize}
        \item Gateway (Stripe, Momo) gửi webhook callback về \texttt{POST /v1/payments/\{id\}/confirm}.
        \item Backend xác thực signature webhook để đảm bảo độ bảo mật.
        \item Cập nhật trạng thái appointment thành ``CONFIRMED''.
    \end{itemize}
    
    \item \textbf{Xác nhận thành công:}
    \begin{itemize}
        \item Gửi push notification cho người dùng: ``Booking confirmed! Your appointment is on [date] at [time] with [specialist]''
        \item Gửi email confirmation với QR code check-in
        \item Gửi notification cho nhà cung cấp và chuyên gia: ``New booking from [customer]''
    \end{itemize}
    
    \item \textbf{Người dùng xem lại:} Lịch hẹn hiện lên trong tab ``Upcoming'' của Appointments screen.
\end{enumerate}

\subsection{Quy trình tư vấn AI}

Quy trình cho AI Health Assistant:

\begin{enumerate}
    \item \textbf{Người dùng mở Chat:} Nhấp vào tab Chat hoặc nút ``AI Health Assistant'' từ Home.
    
    \item \textbf{Khởi tạo conversation:} Nếu chưa có session AI nào, backend tạo một conversation mới.
    \begin{itemize}
        \item Gửi greeting message: ``Hi! I'm Dr. AI 👋''
        \item Hiển thị quick action buttons: Symptom Check, Book Appointment
    \end{itemize}
    
    \item \textbf{Người dùng gửi tin nhắn:} Gõ triệu chứng hoặc câu hỏi (ví dụ: ``I've been having acne for a while'').
    
    \item \textbf{Backend xử lý:}
    \begin{itemize}
        \item Lưu message vào database
        \item Gọi LLM API (OpenAI GPT-4, Anthropic Claude, v.v.) với prompt tinh chỉnh:
        \begin{verbatim}
You are a friendly health advisor. Given the user's symptoms and medical history,
provide helpful advice and recommend relevant medical services from our platform.
Be empathetic and professional. Suggest 1-3 services that match the condition.
        \end{verbatim}
        \item LLM trả về response dạng text + recommendation JSON (với service IDs)
    \end{itemize}
    
    \item \textbf{Streaming Response:} Backend stream từng token của LLM response tới frontend qua Server-Sent Events (SSE) hoặc WebSocket.
    \begin{itemize}
        \item Frontend hiển thị từng từ khi nó đến (typewriter effect)
        \item Giúp tránh cảm giác chờ lâu
    \end{itemize}
    
    \item \textbf{Đề xuất dịch vụ:} Sau lời khuyên chính, AI hiển thị danh sách dịch vụ đề xuất (ví dụ: ``Dermatology Acne Consultation'', ``Basic Facial Care Package'').
    \begin{itemize}
        \item Mỗi dịch vụ có thumbnail, tên, giá, nút ``View →''
        \item Người dùng có thể nhấp để xem chi tiết hoặc bắt đầu đặt lịch
    \end{itemize}
    
    \item \textbf{Tiếp tục hội thoại:} Người dùng có thể hỏi tiếp, AI trả lời dựa trên context cuộc trò chuyện trước đó.
\end{enumerate}

\section{Tối ưu hóa hiệu năng và UX}

\subsection{Caching Strategy}

\begin{itemize}
    \item \textbf{Service Data Cache:} Danh sách dịch vụ được cache cục bộ (SQLite mobile, localStorage web) trong 24 giờ. TTL được set qua HTTP header \texttt{Cache-Control: max-age=86400}.
    
    \item \textbf{User Profile Cache:} Thông tin user được cache 1 giờ. Khi user đăng nhập hoặc update profile, cache được invalidate tức thì.
    
    \item \textbf{Appointment List Cache:} Danh sách appointment được cache 5 phút. Mỗi khi app resumed (bring to foreground), cache được check -- nếu cũ hơn 5 phút, refetch.
    
    \item \textbf{CDN cho Media:} Hình ảnh dịch vụ được lưu trên CDN (CloudFront, Cloudflare), client download từ CDN thay vì origin server.
\end{itemize}

\subsection{Performance Optimization}

\begin{itemize}
    \item \textbf{Image Compression:} Tất cả hình ảnh dịch vụ được compress tới 80\% quality, giảm dung lượng từ 2MB xuống 300-500KB per image.
    
    \item \textbf{Pagination:} Danh sách dịch vụ và appointment được phân trang (per page: 10 items). Người dùng scroll xuống để load batch tiếp theo.
    
    \item \textbf{Lazy Loading:} Component ``Recommend For You'' không được render ngay, chỉ render khi user scroll tới vị trí của nó.
    
    \item \textbf{Code Splitting:} Chat screen, Checkout screen được lazy-load khi user navigate tới, không bundle vào initial app.
    
    \item \textbf{Debouncing:} Khi người dùng tìm kiếm dịch vụ, gọi API search được debounce 300ms để tránh spam request.
\end{itemize}

\subsection{Accessibility}

\begin{itemize}
    \item \textbf{Screen Reader Support:} Tất cả UI element có semantic labels (ánh xạ sang \texttt{semanticLabel} trong Flutter).
    
    \item \textbf{Color Contrast:} Text và background đạt tối thiểu WCAG AA standard (4.5:1 ratio).
    
    \item \textbf{Touch Target Size:} Nút bấm có kích thước tối thiểu 48x48dp, phù hợp với Android Material Design guidelines.
    
    \item \textbf{Dark Mode Support:} Ứng dụng hỗ trợ Dark mode, tự động theo hệ thống hoặc cho phép user chọn trong Settings.
\end{itemize}
